\subsection{Review of Logic}

\content{
        \begin{definition} A \textbf{statement} is either true or false.
        \end{definition}
        \begin{definition} Boolean Logic combines multiple statements
        that are either true or false into an expression that is
        either true or false. 
        \end{definition}
}{% presentation
\vspace{2in}
}{% nopresentation
\vspace{2in}
}{% comments
Mention that three of the most common operations in boolean logic are 'and`,
'or`, and 'not`, and that these correspond to three set operations.
}

\content{
        \begin{definition} A \textbf{conditional} is a compound
        statement of the form "if $p$, then $q$`` or 
        "if $p$, then $q$, else $r$``.
        \end{definition}
}{% presentation
}{% nopresentation
}{% comments
}

\content{
        \begin{definition} A \textbf{universial quantifier} states that an
        entire set shares a characteristic. 
        \end{definition}

        \begin{definition} A \textbf{existential quantifier} states that a set
        contains at least one element.
        \end{definition}
}{% presentation
}{% nopresentation
}{% comments
Something interesting happens when you negate a quantifier...
}

\content{
\begin{definition} We use the following symbols for conjunction (and),
disjunction (or), and negation (not): 
$$ \text{and: } \land,\quad \text{or: }\lor,\quad 
\text{not: }\neg. $$
Sometimes you may see the symbol $\sim$ used for negation instead of $\neg$.
\end{definition}
}{% presentation
}{% nopresentation
}{% comments
}

\content{
\begin{example} Let $C$ be the statement ``I like Coca-Cola" and $P$ be the
statement ``I like Pepsi". Write the following in symbols.
\begin{enumerate}
\item I like Coke and Pepsi
\item I like Coke or I like Pepsi
\item I do not like Coke
\item I like Coke and I do not like Pepsi
\end{enumerate}
\end{example}
}{% presentation
\vspace{1in}
}{% nopresentation
}{% comments
}

\content{
\begin{definition} A \textbf{truth table} is a table showing the value of a
compound statement for all possible values of the simple statements. 
\end{definition}
}{% presentation
}{% nopresentation
}{% comments
}

\content{
\begin{example} Construct a truth table for the compound statement $p\land q$.\\
\begin{tabular}{|c|c|c|}
\hline
$p$ & $q$ &\hspace{0.75cm}  \\
\hline
T & T &  \\
\hline
T & F &  \\
\hline
F & T &  \\
\hline
F & F &  \\
\hline
\end{tabular}
\end{example}
}{% presentation
}{% nopresentation
}{% comments
}

\content{
\begin{example} Construct a truth table for the compound statement 
$\neg (p\land q)$.\\
\begin{tabular}{|c|c|c|c|}
\hline
$p$ & $q$ &\hspace{0.75cm} & \hspace{0.75cm} \\
\hline
T & T & & \\
\hline
T & F & & \\
\hline
F & T & & \\
\hline
F & F & & \\
\hline
\end{tabular}
\end{example}
}{% presentation
}{% nopresentation
}{% comments
}

\content{
\begin{definition}
A conditional is a logical compound statement in which a statement $p$, implies
a statement $q$. This is written $p\Rightarrow q$ (or sometimes $p\rightarrow
q$). The truth table of the conditional $p\Rightarrow q$ is:\\
\begin{tabular}{|c|c|c|}
\hline
$p$ & $q$ & $p \Rightarrow q$ \\
\hline
T & T &  T\\
\hline
T & F &  F\\
\hline
F & T &  T\\
\hline
F & F &  T\\
\hline
\end{tabular}
\end{definition}
}{% presentation
}{% nopresentation
}{% comments
}

\content{
\begin{example} Construct a truth table for the compound statement 
$(p\lor q)\Rightarrow r$.\\
\begin{tabular}{|c|c|c|c|c|c|}
\hline
$p$ & $q$ & $r$ & \hspace{0.75cm} & \hspace{0.75cm} & \hspace{1.5cm} \\
\hline
T & T & T & & & \\
\hline
T & T & F & & &\\
\hline
T & F & T& & & \\
\hline
T & F & F& & &\\
\hline
F & T & T & & &\\
\hline
F & T & F& & & \\
\hline
F & F & T& & &\\
\hline
F & F & F& & &\\
\hline
\end{tabular}
\end{example}
}{% presentation
}{% nopresentation
}{% comments
}

\content{
        \begin{definition} (Analyze Argument with Truth Table) To perform the
        analysis, we do the following steps:
        \begin{enumerate}
        \item Represent each of the premises symbolically
        \item Combine all premises to form a compound statement, this will form
        the antecedent
        \item Create a conditional statement with the conclusion as the
        consequent
        \item Create a truth table for the statement, if it is always true, the
        argument was valid.
        \end{enumerate}
        \end{definition}
}{% presentation
}{% nopresentation
}{% comments
}

\content{
        \begin{example} Consider the argument: 
        \begin{enumerate}
        \item Premise: ``If you bought bread, you went to the store."
        \item Premise: ``You bought bread"
        \item Conclusion: ``You went to the store."
        \end{enumerate}
        \end{example}
}{% presentation
\begin{tabular}{|c|c|c|c|}
\hline
 &  &\hspace{2cm} &\hspace{3cm}  \\
\hline
T & T &&  \\
\hline
T & F &&  \\
\hline
F & T &&  \\
\hline
F & F &&  \\
\hline
\end{tabular}
}{% nopresentation
\begin{tabular}{|c|c|c|c|}
\hline
 &  &\hspace{0.75cm} &\hspace{0.75cm}  \\
\hline
T & T &&  \\
\hline
T & F &&  \\
\hline
F & T &&  \\
\hline
F & F &&  \\
\hline
\end{tabular}
}{% comments
}

\content{
        \begin{example} Consider the argument: 
        \begin{enumerate}
        \item Premise: ``If you listen to the Grateful Dead, you are a hippie."
        \item Premise: ``Sky does not listen to the Grateful Dead"
        \item Conclusion: ``Sky is not a hippie."
        \end{enumerate}
        \end{example}
}{% presentation
\begin{tabular}{|c|c|c|c|c|}
\hline
 &  &\hspace{0.75cm} &\hspace{0.75cm} & \hspace{2cm}  \\
\hline
T & T &&&  \\
\hline
T & F &&&  \\
\hline
F & T &&&  \\
\hline
F & F &&&  \\
\hline
\end{tabular}
}{% nopresentation
\begin{tabular}{|c|c|c|c|}
\hline
 &  &\hspace{0.75cm} &\hspace{0.75cm}  \\
\hline
T & T &&  \\
\hline
T & F &&  \\
\hline
F & T &&  \\
\hline
F & F &&  \\
\hline
\end{tabular}
}{% comments
}
